\section{Introduction}
\label{sec:introduction}

The advantages of modeling data with algebraic datatypes (ADTs) 
have been recognised since the 80s 
\cite{Burstall1977a,DBLP:conf/lfp/BurstallMS80}. 
% ADTs are easy to define, 
% promote abstraction, and are less prone to errors.
ADTs are easy to define and use, they enable the construction of
domain-specific types, and are amenable to efficient nominal type checking.
Functional programming languages quickly embraced ADTs along
with the concept of pattern matching which is still gaining momentum.  

Algebraic datatypes integrate recursive types, sum types, product
types, and generativity (each ADT definition creates a new type) in a
single declaration. The smooth integration of different typing concepts
makes ADTs easy to learn and use. It also simplifies the 
implementation of type checkers because they can elide superlinear algorithms for checking
equivalence of equirecursive types (where recursive types are
considered equal to their unfolding) in favor of a linear comparison of nominal types.
%  and efficient to handle in a language
% implementation. ADTs are popular among programmers as they have to learn just a single mechanism
% for type definition.

We show that the ADT approach for defining recursive
datatypes can be carried over to 
\emph{algebraic protocols} (AP), a novel method for defining a large class
of recursive protocols. Analogous to ADTs, APs integrate 
concepts from the theory of session types
\cite{DBLP:conf/concur/Honda93,DBLP:conf/parle/TakeuchiHK94}
(recursive types, choice types, and sequence types) and generativity
in a single declaration. Analogous to ADTs, APs are easy to use and
are amenable to efficient type checking. 

We develop APs by analogy to ADTs. Consider the following ADT for lists of
integers.
\begin{lstlisting}
  data IntList = Nil | Cons Int IntList
\end{lstlisting}
The datatype \lstinline+IntList+  comes with two data constructors,
\lstinline+Nil+  to construct an 
empty list and \lstinline+Cons+ to construct a list from an element
of type \lstinline+Int+ and a list of type \lstinline+IntList+. The
alternative between the constructors corresponds to a sum type and the
two arguments of \lstinline+Cons+ give rise to a product
type. Clearly, \lstinline+IntList+ is a recursive type. Moreover, its
definition introduces \lstinline+IntList+ as a new type, different
from any other. Type equivalence is simple: we just compare ADTs
by name.

Although \lstinline+IntList+ specifies a (static) data type, 
it carries an implicit dynamics and it should be 
quite natural for the reader to imagine elements of this type as
finite sequences of messages between two communication partners: if
there are  further elements, the sender selects the  
constructor \lstinline+Cons+, sends a value of type 
\lstinline+Int+ and then recurses. If there are no further elements,
the sender selects \lstinline+Nil+ and stops the transmission. 
On the other end, the receiver awaits a constructor. On receiving a
\lstinline+Cons+ it receives an \lstinline+Int+ and recurses. It stops
on receiving a \lstinline+Nil+. 

% perhaps, we could rather \emph{receive} a value of type \lstinline+a+ 
% and then receive the rest of the stream. In this paper we 
% interpret algebraic datatypes as protocols and reformulate 
% the theory of session types, now presented as a natural 
% extension of algebraic datatypes -- materialised in our
% theory of \emph{algebraic session types}. 

This view has striking parallels to the theory of binary session
types~\cite{DBLP:conf/concur/Honda93,DBLP:conf/esop/HondaVK98,DBLP:conf/parle/TakeuchiHK94}.
Session types provide a rich type discipline for communication channels that
statically guarantees protocol fidelity as well as type soundness across
communication channels. Here is a session type for the protocol associated to
the \lstinline+IntList+ ADT, written in the syntax of existing
systems~\cite{DBLP:journals/corr/abs-2106-06658,lindley17:_light_funct_session_types,DBLP:journals/jfp/GayV10}
and as seen from the sender's perspective:\footnote{Writing
  \lstinline+IntListS+ for the session type corresponding to
  \lstinline+IntList+.}
%
\begin{lstlisting}
  type IntListS = mualpha. oplus{Cons: !Int.alpha, Nil: EndT}
\end{lstlisting}
%
The leading $\mu\alpha\dots$ constructs an equirecursive type. The body is
an internal choice type $\oplus$ with two alternatives indicated by the tags
\lstinline+Cons+ and \lstinline+Nil+. If the sender selects
\lstinline+Cons+, it must continue according to
\lstinline+!Int.alpha+, which means it must send an integer and then
recurse. If the sender selects \lstinline+Nil+, then the protocol
comes to an end as indicated by the type \lstinline+EndT+.
%
Type \lstinline+EndT+ represents a channel endpoint ready to be closed. The dual
type, \lstinline+EndW+, represents the other endpoint, waiting to be closed.
%
As the name
suggests, a session type describes a full session of interactions
between the communication partners up to the end.

Like this example, most session type systems 
rely on structural typing and equirecursion; 
this way, unfolding does not give rise to communication.
Traditional session type systems restrict 
recursion to the tail position (as in \lstinline/IntListS/), so that types are effectively finite
automata with regular trace languages. Consequently, type equivalence
and subtyping for such systems are decidable in polynomial time.  Recently, richer forms of recursion
have been considered that lift the restriction to tail recursion, so
that types are akin to, e.g.,  deterministic context-free grammars
\cite{DBLP:conf/icfp/ThiemannV16,DBLP:conf/esop/Padovani17,DBLP:journals/toplas/Padovani19,DBLP:journals/corr/abs-2106-06658,
  DBLP:conf/esop/DasDMP21}.  
Unfortunately, when types are endowed with richer forms 
of recursion, type equivalence gets more complex 
(doubly-exponential \cite{DBLP:conf/esop/DasDMP21}). 
The implementation of a type equivalence algorithm for these systems is nontrivial
\cite{DBLP:conf/tacas/AlmeidaMV20}; incomplete algorithms have been
used to speed up type equivalence \cite{DBLP:conf/esop/DasDMP21};
annotations have been proposed~\cite{DBLP:journals/toplas/Padovani19}
that make type checking effective. But the algorithms for
type equivalence remain asymptotically superlinear in all these works.
% type checking by exploiting programmer annotations and
% showed that subtyping is undecidable. 

The sad truth is that structural session types are quite verbose and
too complex for the average programmer, thus hampering their adoption. While equirecursive types are
expressive, understanding their equivalence is hard for programmers
and inefficient in language implementations.

What if programming with expressive recursive protocols were as easy as
programming with ADTs? The theory and practice of algebraic protocols
that we propose in this paper delivers exactly such an
approach. Enhanced with parameters and polymorphism in the style of
System~F, APs provide a uniquely simple and efficient approach to modularity and abstraction
when defining and implementing communication protocols.

%Continuing the \lstinline+IntList+ example,
Here is the declaration of an algebraic
protocol corresponding to the session type \lstinline+IntListS+:\footnote{Writing
  \lstinline+IntListP+ for the protocol type corresponding to
  \lstinline+IntList+.}
\begin{lstlisting}
  protocol IntListP = Nil | Cons Int IntListP
\end{lstlisting}
% In this paper, we propose a paradigm shift. We revisit the 
% guidelines that constituted the basis of the functional languages 
% that we use today 
% \cite{Burstall1977a,DBLP:conf/lfp/BurstallMS80}
% and reformulate the theory of session types 
% by favouring modularity, code abstraction, the extensive use of 
% user-defined types and overloaded operators.
% Algebraic session types now enable the definition of protocol 
% templates that are \emph{not} committed to a direction (of communication)
% until they are \emph{materialized}. The protocol template for 
% streams reads like:
% \begin{lstlisting}
%     protocol StreamP a = Cons a (StreamP a)
% \end{lstlisting}
The keyword \lstinline+protocol+, in place of \lstinline+data+, denotes an
algebraic \emph{protocol} type rather than an algebraic \emph{data} type. A
protocol type is like a template that describes a composable segment of a
session. It neither describes a full session, nor is it committed to a
particular direction of communication. Rather, a protocol type has to be
\emph{materialized} into a session type by assigning a direction and specifying
a continuation session: we write \lstinline+!IntListP.S+ to send a sequence of
integers and continue with session type \lstinline+S+ or \lstinline+?IntListP.S+
to receive a sequence of integers and continue with \lstinline+S+, so that the
session type \lstinline/IntListS/ corresponds to \lstinline/!IntListP.EndT/.
Hence, protocol types relate to session types just like composable continuations
relate to continuations \cite{DBLP:conf/popl/Felleisen88,DanvyFilinski1992}: the
former describe a segment of interaction (execution) while the latter describe
interaction (execution) to the end.

Just like ADTs, we can equip APs with parameters as in
\lstinline+ListP a+:
\begin{lstlisting}
  protocol ListP a = Nil | Cons a ListP
\end{lstlisting}
While parameters range over types in ADT definitions, they
range over \emph{protocols} in AP definitions! This facility is key to
the modularity of APs. For example, the protocol
\lstinline+ListP Q+ repeats the protocol \lstinline+Q+. It is
also possible to flip the direction of \lstinline+Q+ by considering the protocol
\lstinline+ListP -Q+: repeat the protocol \lstinline+Q+ but with reversed roles
of sender and receiver (more in \cref{sec:simple-protocols-add}). With this
modular design, we can define finite streams of
anything expressible as an algebraic protocol. The
external and internal choices in previous session type systems are
captured by the sum-of-product nature of algebraic types---the sum
type replaced by internal/external choice and the product type
replaced by sequencing as in the linear logic
interpretation of session types \cite{DBLP:journals/mscs/CairesPT16}.

Algebraic protocol types are nominal and isorecursive. 
This combination of features reduces the complexity of type equivalence from
asymptotically doubly-exponential \cite{DBLP:conf/tacas/AlmeidaMV20,DBLP:conf/esop/DasDMP21} 
to linear. Notably, the expressivity of algebraic protocols 
is on par with that of session types with non-regular recursion
\cite{DBLP:conf/icfp/ThiemannV16,DBLP:conf/esop/DasDMP21}, 
so all protocols expressible in such theories can be adapted 
to algebraic protocols (without paying the price of hosting complex or 
incomplete algorithms in the compiler). Interestingly, types 
are now also iso-associative, thus saving us from defining 
hard-coded associativity laws, due to the basic properties 
of sequential composition 
\cite{DBLP:conf/icfp/ThiemannV16,DBLP:journals/corr/abs-2106-06658}.

\subsection*{Contributions}
\label{sec:contributions}

\begin{itemize}
\item Algebraic protocols (AP). We propose a new approach for defining
  recursive protocols compositionally inspired by parameterized algebraic
  datatypes. This approach combines naming, recursion, and
  choice-of-sequence types analogously to algebraic
  datatypes. It adds a novel notion of direction to the component
  protocols of a tag. Checking type equivalence becomes simple and efficient
  (worst case linear in the size of the types) as APs are nominal. 
\item Algebraic session types (\algst). We define a core calculus that is based
  on linear System~F with (non-linear) recursion and features session types
  based on APs. Type checking for \algst is specified by bidirectional typing
  rules that have a direct algorithmic reading.
  % 
  % The operational semantics is presented as a labelled transition relation,
  % inspired by Fowler \etal and Montesi and
  % Peressotti~\cite{DBLP:conf/concur/0001KDLM21,DBLP:journals/corr/abs-2105-08996,DBLP:journals/corr/abs-2106-11818}
  % for double binders and the open/close scheme for the free output of the
  % $\pi$-calculus. %~\cite{DBLP:journals/iandc/MilnerPW92b}.
  We prove type soundness via preservation and
  progress.% , where our approach offers a
  % rather elegant statement (\cref{thm:progress} on \cpageref{thm:progress}).
\item Full implementation as a publicly available artifact
  \cite{artifact}, including practical extensions (cf.\ \cref{sec:implementation}).
\end{itemize}

\subsection*{Overview}
\label{sec:overview}

\Cref{sec:algebr-sess-types} contains an informal introduction of algebraic
protocols with examples; we illustrate the gain in
expressivity and modularity with respect to other work in
\cref{sec:param-prot}. \Cref{sec:types} discusses the type structure
of \algst and establishes key results for efficient type checking: correctness of type normalization and
worst case complexity of type equivalence. \Cref{sec:processes} defines
statics and dynamics of expressions 
and processes % . \Cref{sec:metatheory}
followed by a presentation of the metatheoretical results.
% \Cref{sec:cfst-vs-algst} gives an informal outline of
% an embedding of context-free session types to \algst.
\Cref{sec:implementation} discusses the implementation; since all rules are
algorithmic, their appropriate reading leads directly to a type checker and an interpreter.
\Cref{sec:related,sec:conclusion} discuss related work and conclude.
%
Proofs, auxiliary results, further examples, an informal outline of an embedding
of context-free session types to \algst, and further discussion are provided as supplementary material. The
implementation is publicly available \cite{artifact}.

%%% Local Variables:
%%% mode: latex
%%% TeX-master: "main"
%%% End:
